% Bagian: first-order-logic
% Bab: axiomatic-deduction
% Subbagian: identity

\documentclass[../../../include/open-logic-section]{subfiles}

\begin{document}

\olfileid[id]{fol}{axd}{ide}

\olsection{\usetoken{P}{derivation} dengan \usetoken{S}{identity}}

Untuk mengakomodasi $\eq$ dalam !!{derivation}s, kita cukup menambahkan
skema aksioma baru. Definisi !!{derivation} dan $\Proves$ tetap sama;
kita hanya mengizinkan aksioma-aksioma baru itu juga.

\begin{defn}[Aksioma untuk !!{identity}]
\begin{align}
\ollabel{ax:id1} & \eq[t][t], \\
\ollabel{ax:id2} & \eq[t_1][t_2] \lif (!B(t_1) \lif
  !B(t_2)),
\end{align}
untuk sebarang suku tertutup $t$, $t_1$, dan $t_2$.
\end{defn}

\begin{prop}\ollabel{prop:sound}
  Aksioma \olref{ax:id1} dan \olref{ax:id2} bersifat valid.
\end{prop}

\begin{proof}
  Latihan.
\end{proof}

\begin{prob}
Buktikan \olref[fol][axd][ide]{prop:sound}.
\end{prob}

\begin{prop}\ollabel{prop:iden1}
 $\Gamma \Proves \eq[t][t]$ untuk sebarang suku tertutup $t$ dan
 himpunan~$\Gamma$.
\end{prop}

\begin{prop}\ollabel{prop:iden2}
  Untuk suku-suku tertutup $t_1$ dan $t_2$, jika
  $\Gamma \Proves !A(t_1)$ dan $\Gamma \Proves
  \eq[t_1][t_2]$, maka $\Gamma \Proves !A(t_2)$.
\end{prop}

\begin{proof}
!!^{formula}
\[
(\eq[t_1][t_2] \lif (!A(t_1) \lif !A(t_2)))
\]
merupakan instans~\olref{ax:id2}. Kesimpulannya diperoleh melalui dua
penerapan~\MP.
\end{proof}

\end{document}
